% Part: first-order-logic
% Chapter: axiomatic-deduction
% Section: rules-and-proofs

\documentclass[../../../include/open-logic-section]{subfiles}

\begin{document}

\iftag{FOL}
      {\olfileid{fol}{axd}{rul}}
      {\olfileid{pl}{axd}{rul}}

\olsection{规则与推导}

\begin{explain}
  公理化推导或许是逻辑学中最简单的推导系统。推导只是一个公式序列。但要算作推导，序列中的每个公式要么是某条公理的实例，要么是应用某条推理规则，从序列中排在它前面的一个或多个公式推出的。推导得出的是其最后一个公式。
\end{explain}

\begin{defn}[推导]
如果 $\Gamma$ 是语言 $\Lang L$ 的一个公式集，那么从~$\Gamma$ 出发的一个\emph{推导}就是一个公式的有限序列 $!A_1$, \dots, $!A_n$，其中对每个 $i \le n$，以下条件之一成立：
\begin{enumerate}
\item $!A_i \in \Gamma$；或者
\item $!A_i$ 是一条公理；或者
\item $!A_i$ 是应用某条推理规则，由满足 $j < i$ 的某个 $!A_j$（以及满足 $k < i$ 的某个 $!A_k$）推出的。
\end{enumerate}
\end{defn}

什么算作正确的推导，取决于我们允许哪些推理规则（当然，也取决于我们把什么当作公理）。而推理规则是一个“如果–那么”陈述，它告诉我们在某些条件下，推导中某一步~$!A_i$ 是一个正确的推理步骤。

\begin{defn}[推理规则]
一条\emph{推理规则}给出了从~$\Gamma$ 出发的推导中，什么算作正确推理步骤的一个充分条件。
\end{defn}

例如，只要 $!A \in \Gamma$，任何单元素序列 $!A$ 都平凡地算作一个推导，因此下面这条可以是一条非常简单的推理规则：

\begin{quote}
  若 $!A \in \Gamma$，则 $!A$ 在任何从~$\Gamma$ 出发的推导中总是一个正确的推理步骤。
\end{quote}

类似地，如果 $!A$ 是一条公理，那么 $!A$ 本身就是一个推导，因此这也是一条推理规则：

\begin{quote}
  若 $!A$ 是一条公理，则 $!A$ 是一个正确的推理步骤。
\end{quote}

当推理规则依赖于所考虑步骤之前出现的公式时，事情就变得更有趣了。下面这条规则被称为\emph{肯定前件式}：

\begin{quote}
  若 $!B \lif !A$ 和 $!B$ 出现在推导中更高处（即更早），则 $!A$ 是一个正确的推理步骤。
\end{quote}

如果这是唯一的推理规则，那么我们上面给出的推导定义就等价于：$!A_1$, \dots, $!A_n$ 是一个推导，当且仅当对每个 $i \le n$，以下条件之一成立：

\begin{enumerate}
\item $!A_i \in \Gamma$；或者
\item $!A_i$ 是一条公理；或者
\item 对某个 $j < i$，$!A_j$ 是 $!B \lif !A_i$，并且对某个 $k < i$，$!A_k$ 是 $!B$。
\end{enumerate}

最后一个条款说的是：$!A_i$ 由 $!A_j$（即 $!B \lif !A_i$）和 $!A_k$（即 $!B$）通过肯定前件式得到。如果我们从 $1$ 逐步检查到 $n$，每次都发现公式 $!A_i$ 要么属于 $\Gamma$，要么是公理，要么根据某条推理规则是一个正确的推理步骤，那么整个序列就可以算作一个正确的推导。

\begin{defn}[可推导性]
公式~$!A$ 是\emph{可从 $\Gamma$ 推导的}，记作 $\Gamma \Proves !A$，如果存在一个从 $\Gamma$ 出发、以 $!A$ 结尾的推导。
\end{defn}

\begin{defn}[定理]
公式~$!A$ 是一个\emph{定理}，如果存在从空集出发到 $!A$ 的推导。若 $!A$ 是定理，则记作 $\Proves !A$；否则记作 $\Proves/ !A$。
\end{defn}

\end{document}